% =============================================================================
% SECCION: VALIDACION ESTADISTICA
% =============================================================================

\subsection{Validaci\'on Estad\'istica}
\label{sec:statistical_validation}

M\'as all\'a de la comparaci\'on directa de retornos presentada en secciones anteriores, se requiere validaci\'on estad\'istica formal para evaluar en qu\'e medida los resultados observados podr\'ian ser atribuibles al azar muestral. A continuaci\'on se presentan los tests estad\'isticos aplicados a los cuatro modelos ganadores (LSTM+Attention, Ridge, RandomForest y CNN-LSTM) y al \textit{benchmark} SPY \textit{Buy \& Hold}, todos computados a partir de los retornos diarios del per\'iodo de prueba (1{,}302 d\'ias).

% ─────────────────────────────────────────────────────────────────────────────
\subsubsection{Bootstrap de Sharpe Ratios}
\label{subsec:bootstrap_sharpe}

Se emplea \hyperlink{glos:bootstrap}{\textit{bootstrap}} no param\'etrico \citep{efron1979bootstrap} con $B = 5{,}000$ muestras con reemplazo para construir intervalos de confianza al 95\% del Sharpe Ratio de cada estrategia. Se incorpora la correcci\'on de \citet{lo2002} al error est\'andar para ajustar por la autocorrelaci\'on de primer orden de los retornos.

\begin{equation}
\text{SE}(\hat{S}) = \sqrt{\frac{1 + \frac{\hat{S}^2}{2} \cdot (1 + \hat{\rho}_1)}{T \cdot (1 - \hat{\rho}_1)}}
\label{eq:lo2002}
\end{equation}

Los resultados (Tabla~\ref{tab:sharpe_ci}) sugieren que LSTM+Attention es el \'unico modelo cuyo intervalo de confianza se mantiene \'integramente positivo (0.375--2.424), con 99.6\% de probabilidad de Sharpe positivo y 83.1\% de superar al \textit{benchmark}. Ridge alcanza 94.9\% de probabilidad positiva pero su l\'imite inferior se mantiene ligeramente negativo ($-$0.121), mientras que RandomForest y CNN-LSTM presentan probabilidades cercanas al 50\% de superar al SPY (49.9\% y 44.4\%), indicando evidencia d\'ebil de que el mejor rendimiento observado sea estad\'isticamente distinguible del \textit{benchmark} pese a retornos absolutos positivos.

% ─────────────────────────────────────────────────────────────────────────────
\subsubsection{Test de Diebold-Mariano}
\label{subsec:diebold_mariano}

El test de \citet{diebold1995} eval\'ua si la diferencia en rendimiento econ\'omico entre dos estrategias es estad\'isticamente significativa, definiendo $d_t = L(\text{benchmark}_t) - L(\text{estrategia}_t)$ con funci\'on de p\'erdida igual al retorno negativo.

\begin{equation}
DM = \frac{\bar{d}}{\sqrt{\hat{V}(\bar{d})}}
\label{eq:dm_test}
\end{equation}

donde $\hat{V}(\bar{d})$ incorpora la correcci\'on de \citet{newey1987} con cinco rezagos. La Tabla~\ref{tab:dm_test} sugiere que LSTM+Attention es el \'unico modelo que alcanza significancia al 5\% ($DM = 2.296$, $p = 0.011$), permitiendo rechazar la hip\'otesis nula de retornos indistinguibles del \textit{benchmark}. Ridge se sit\'ua en la zona marginal ($p = 0.101$), mientras que RandomForest y CNN-LSTM no presentan evidencia de mejor rendimiento respecto al \textit{benchmark} ($p > 0.35$), lo que ilustra la asimetr\'ia habitual entre significancia econ\'omica y estad\'istica en la predicci\'on de retornos.

% ─────────────────────────────────────────────────────────────────────────────
\subsubsection{Test de Clark-West}
\label{subsec:clark_west}

El test de \citet{clark2007approximately} eval\'ua la capacidad predictiva de modelos anidados comparando cada modelo contra el predictor de media hist\'orica ($\bar{r}_{\text{train}} = 0.025\%$ diario).

\begin{equation}
CW_t = (e_{1,t}^2 - e_{2,t}^2) + (\hat{y}_{1,t} - \hat{y}_{2,t})^2
\label{eq:cw_test}
\end{equation}

La Tabla~\ref{tab:cw_test} es consistente con el patr\'on identificado, dado que LSTM+Attention es nuevamente el \'unico modelo significativo ($CW = 1.901$, $p = 0.029$). La convergencia de ambos tests, que eval\'uan dimensiones distintas del rendimiento (retornos de estrategia vs.\ error de predicci\'on), sobre el mismo modelo es compatible con la interpretaci\'on de un error de predicci\'on menor al del predictor de media hist\'orica en este per\'iodo. El valor negativo de RandomForest ($CW = -0.336$) indica que sus predicciones no mejoran respecto a la media hist\'orica, sugiriendo que la rentabilidad de la estrategia proviene de se\~nales direccionales asim\'etricas \citep{pesaran2002}.

% ─────────────────────────────────────────────────────────────────────────────
\subsubsection{Model Confidence Set}
\label{subsec:mcs}

El \textit{Model Confidence Set} (MCS) de \citet{hansen2011} identifica el subconjunto de modelos que contiene al mejor con probabilidad dada, mediante el estad\'istico $T_{\max}$ con \textit{bootstrap} de bloques estacionarios ($B = 3{,}000$) al nivel $\alpha = 0.10$. Aplicado a las 24 estrategias (23 modelos + SPY), el procedimiento no logra eliminar ning\'un modelo del conjunto.

\begin{equation}
\hat{\mathcal{M}}^*_{0.10} = \mathcal{M}_0 \quad \text{(24 modelos, 0 eliminados)}
\label{eq:mcs_result}
\end{equation}

Este resultado se explica por la elevada varianza de los retornos diarios de renta variable, dado que la desviaci\'on est\'andar diaria del S\&P~500 ($\approx 1.1\%$) excede ampliamente la diferencia media entre estrategias ($\approx 0.05\%$), generando una relaci\'on se\~nal-ruido insuficiente para discriminar con 1{,}302 observaciones, resultado consistente con \citet{hansen2011}.

% ─────────────────────────────────────────────────────────────────────────────
\subsubsection{An\'alisis de \textit{Overfitting} y Comparaciones M\'ultiples}
\label{subsec:overfitting}

La Tabla~\ref{tab:overfitting} muestra que LSTM+Attention exhibe la menor degradaci\'on en precisi\'on direccional entre entrenamiento y prueba ($-$4.6\%), con un error cuadr\'atico que \textit{disminuye} en el per\'iodo fuera de muestra ($0.68\times$), lo que sugiere que el mecanismo de atenci\'on captura patrones con cierta capacidad de generalizaci\'on. Ridge y RandomForest muestran mayor degradaci\'on ($-$18.6\% y $-$20.9\%) con ratios de MSE superiores a uno, indicando sobreajuste m\'as pronunciado aunque sin anular la rentabilidad. XGBoost, como contraejemplo, muestra m\'inima degradaci\'on ($-$1.6\%) pero un Sharpe de $-$0.024, dado que su baja degradaci\'on refleja ausencia de se\~nal en ambos per\'iodos.

Con 23 modelos evaluados simult\'aneamente, la correcci\'on de Bonferroni establece un umbral de $\alpha_{\text{adj}} = 0.05/23 \approx 0.0022$ bajo el cual ning\'un modelo alcanza significancia ($p_{\min} = 0.011$ para LSTM+Attention). El procedimiento de \citet{benjamini1995controlling} tampoco identifica modelos significativos al nivel $q = 0.10$, lo que indica que, bajo correcciones est\'andar por comparaciones m\'ultiples, la evidencia de mejor rendimiento individual no alcanza umbrales convencionales.

% ─────────────────────────────────────────────────────────────────────────────
\subsubsection{S\'intesis de la Validaci\'on Estad\'istica}
\label{subsec:validation_summary}

\begin{table}[htbp]
\centering
\caption{Resumen de tests estad\'isticos por modelo}
\label{tab:validation_summary}
\small
\begin{tabular}{@{}lccccc@{}}
\toprule
\textbf{Modelo} & \textbf{P(S$>$0)} & \textbf{P(S$>$SPY)} & \textbf{DM (5\%)} & \textbf{CW (5\%)} & \textbf{Bonferroni} \\
\midrule
LSTM+Attention & 99.6\% & 83.1\% & \cmark & \cmark & $\times$ \\
Ridge & 94.9\% & 60.4\% & $\times$ & $\times$ & $\times$ \\
RandomForest & 93.3\% & 49.9\% & $\times$ & $\times$ & $\times$ \\
CNN-LSTM & 89.3\% & 44.4\% & $\times$ & $\times$ & $\times$ \\
\bottomrule
\end{tabular}
\vspace{0.2cm}
\parbox{0.9\textwidth}{\footnotesize\textit{\cmark{} = significativo al nivel indicado. $\times$ = no significativo. El MCS no se incluye porque no logr\'o eliminar ning\'un modelo del conjunto completo.}}
\end{table}

La Tabla~\ref{tab:validation_summary} sintetiza los resultados. LSTM+Attention es el \'unico modelo que alcanza significancia en los tests de Diebold-Mariano y Clark-West al 5\%, mantiene un intervalo \textit{bootstrap} \'integramente positivo y presenta la menor degradaci\'on predictiva, y la coincidencia de varios tests que eval\'uan dimensiones distintas del rendimiento parece indicar que el mejor rendimiento observado en esta muestra no ser\'ia enteramente atribuible al azar.

No obstante, tres resultados atemperan esta conclusi\'on. Bajo Bonferroni, incluso LSTM+Attention pierde significancia; el MCS no discrimina entre las 24 estrategias dado que la varianza diaria genera una relaci\'on se\~nal-ruido insuficiente con 1{,}302 observaciones; y los tres modelos restantes no alcanzan significancia individual al 5\%. Estos hallazgos matizan la conclusi\'on de \citet{gu2020} sobre la superioridad general de m\'etodos no lineales, puesto que en este contexto solo el modelo con mecanismo de atenci\'on alcanza los umbrales de significancia aplicados, hallazgo consistente con la advertencia de \citet{lopezdeprado2018} sobre los riesgos de evaluar m\'ultiples estrategias simult\'aneamente.
